% GENERATED FILE. Edit canonical lessons, then run npm run generate:books.
% canonical-source-hash: fnv1a64:5a93674e148ea7be
% canonical-lessons: LA-C544-cithara, LA-C544-citharoedus, LA-C544-chorda, LA-C544-fidicen, LA-C544-psallo, LA-R544-review-537-544

\chapter{Strings}
\label{ch:la-strings}
\hlchaptermodality{544}

\begin{chapteropening}
\textbf{By the end of this chapter:} \emph{I can talk about the lyre and those who play it.}

Complete the last lesson of chapter 544: I can talk about the lyre and those who play it.
\end{chapteropening}

\section[cithara, citharae]{cithara, citharae — a cithara, concert lyre}

\label{lesson:LA-C544-cithara}



\begin{quote}
Before the new one: say the Latin for to sing together, then the Latin for melodious.
\end{quote}

\subsection*{You'll want to know: cithara, citharae}
\textbf{cithara, citharae} — \textquotedblleft{}a cithara, concert lyre\textquotedblright{}.

\begin{quote}
Iuvenis citharam pulsat et canit.
\end{quote}

\begin{quote}
\emph{The young man plays the lyre and sings.}
\end{quote}

Feminine, first declension. A large lyre for public performance.

\begin{cousinweb}
\textbf{cithara} was borrowed from Greek \emph{kithara} (Lewis \& Short; EtymDB). English \emph{guitar} and \emph{zither} come from the same Greek word by other routes.
\end{cousinweb}

\begin{grammarlens}[title={What you've built}]
One more word for telling, explaining and describing at B1.
\end{grammarlens}

\subsection*{Your turn}
\begin{itemize}
\raggedright
  \item \emph{Say it:} \emph{cithara, citharae}
  \item \emph{Say it:} \emph{cithara, citharae}, once more
  \item \emph{Say it:} a sentence of your own with \emph{cithara, citharae}
  \item \emph{Recall:} say the Latin for to sing together, then the Latin for melodious, then say \emph{cithara, citharae} again
\end{itemize}

\begin{tcolorbox}[breakable,colback=teal!4,colframe=teal!35!black,title={Before you move on}]
What does cithara, citharae mean? (\textquotedblleft{}A cithara, concert lyre\textquotedblright{}.) Say it once more.
\end{tcolorbox}
\section[citharoedus, citharoedī]{citharoedus, citharoedī — a singer to the lyre}

\label{lesson:LA-C544-citharoedus}



\begin{quote}
Before the new one: say the Latin for melodious, then the Latin for a cithara.
\end{quote}

\subsection*{You'll want to know: citharoedus, citharoedī}
\textbf{citharoedus, citharoedī} — \textquotedblleft{}a singer to the lyre\textquotedblright{}.

\begin{quote}
Citharoedus in theātrō canit.
\end{quote}

\begin{quote}
\emph{The lyre-singer sings in the theatre.}
\end{quote}

Masculine, second declension. A star performer who played the cithara and sang.

\begin{cousinweb}
\textbf{citharoedus} was borrowed from Greek \emph{kitharodos}, built on \emph{kithara} \textquotedblleft{}lyre\textquotedblright{} and \emph{aoidos} \textquotedblleft{}singer\textquotedblright{} (EtymDB; Lewis \& Short).
\end{cousinweb}

\begin{grammarlens}[title={What you've built}]
One more word for telling, explaining and describing at B1.
\end{grammarlens}

\subsection*{Your turn}
\begin{itemize}
\raggedright
  \item \emph{Say it:} \emph{citharoedus, citharoedī}
  \item \emph{Say it:} \emph{citharoedus, citharoedī}, once more
  \item \emph{Say it:} a sentence of your own with \emph{citharoedus, citharoedī}
  \item \emph{Recall:} say the Latin for melodious, then the Latin for a cithara, then say \emph{citharoedus, citharoedī} again
\end{itemize}

\begin{tcolorbox}[breakable,colback=teal!4,colframe=teal!35!black,title={Before you move on}]
What does citharoedus, citharoedī mean? (\textquotedblleft{}A singer to the lyre\textquotedblright{}.) Say it once more.
\end{tcolorbox}
\section[chorda, chordae]{chorda, chordae — a string (of an instrument)}

\label{lesson:LA-C544-chorda}



\begin{quote}
Before the new one: say the Latin for a cithara, then the Latin for a singer to the lyre.
\end{quote}

\subsection*{You'll want to know: chorda, chordae}
\textbf{chorda, chordae} — \textquotedblleft{}a string (of an instrument)\textquotedblright{}.

\begin{quote}
Ūna chorda lyrae rupta est.
\end{quote}

\begin{quote}
\emph{One string of the lyre is broken.}
\end{quote}

Feminine, first declension, from Greek: first gut, then a string. English \emph{chord} and \emph{cord}.

\begin{cousinweb}
\textbf{chorda} was borrowed from Greek \emph{chorde} \textquotedblleft{}gut, string\textquotedblright{} (Lewis \& Short; EtymDB). English \emph{cord} and \emph{chord} come from it.
\end{cousinweb}

\begin{grammarlens}[title={What you've built}]
One more word for telling, explaining and describing at B1.
\end{grammarlens}

\subsection*{Your turn}
\begin{itemize}
\raggedright
  \item \emph{Say it:} \emph{chorda, chordae}
  \item \emph{Say it:} \emph{chorda, chordae}, once more
  \item \emph{Say it:} a sentence of your own with \emph{chorda, chordae}
  \item \emph{Recall:} say the Latin for a cithara, then the Latin for a singer to the lyre, then say \emph{chorda, chordae} again
\end{itemize}

\begin{tcolorbox}[breakable,colback=teal!4,colframe=teal!35!black,title={Before you move on}]
What does chorda, chordae mean? (\textquotedblleft{}A string (of an instrument)\textquotedblright{}.) Say it once more.
\end{tcolorbox}
\section[fidicen, fidicinis]{fidicen, fidicinis — a lyre-player}

\label{lesson:LA-C544-fidicen}



\begin{quote}
Before the new one: say the Latin for a singer to the lyre, then the Latin for a string.
\end{quote}

\subsection*{You'll want to know: fidicen, fidicinis}
\textbf{fidicen, fidicinis} — \textquotedblleft{}a lyre-player\textquotedblright{}.

\begin{quote}
Fidicen ad convīvium venit.
\end{quote}

\begin{quote}
\emph{A lyre-player comes to the dinner party.}
\end{quote}

Masculine, third declension: \emph{fidēs} 'lyre strings' + \emph{canō}.

\begin{cousinweb}
\textbf{fidicen} is built on \emph{fidēs} \textquotedblleft{}the strings of a lyre\textquotedblright{} and \emph{canō} \textquotedblleft{}to sing, play\textquotedblright{} (Lewis \& Short).
\end{cousinweb}

\begin{grammarlens}[title={What you've built}]
One more word for telling, explaining and describing at B1.
\end{grammarlens}

\subsection*{Your turn}
\begin{itemize}
\raggedright
  \item \emph{Say it:} \emph{fidicen, fidicinis}
  \item \emph{Say it:} \emph{fidicen, fidicinis}, once more
  \item \emph{Say it:} a sentence of your own with \emph{fidicen, fidicinis}
  \item \emph{Recall:} say the Latin for a singer to the lyre, then the Latin for a string, then say \emph{fidicen, fidicinis} again
\end{itemize}

\begin{tcolorbox}[breakable,colback=teal!4,colframe=teal!35!black,title={Before you move on}]
What does fidicen, fidicinis mean? (\textquotedblleft{}A lyre-player\textquotedblright{}.) Say it once more.
\end{tcolorbox}
\section[psallō, psallere]{psallō, psallere — to play the cithara, sing to strings}

\label{lesson:LA-C544-psallo}



\begin{quote}
Before the new one: say the Latin for a string, then the Latin for a lyre-player.
\end{quote}

\subsection*{You'll want to know: psallō, psallere}
\textbf{psallō, psallere} — \textquotedblleft{}to play the cithara, sing to strings\textquotedblright{}.

\begin{quote}
Puella psallit et saltat.
\end{quote}

\begin{quote}
\emph{The girl plays the lyre and dances.}
\end{quote}

\emph{psallō, psallere, psallī}. From Greek. English \emph{psalm} comes from the same Greek verb.

\begin{cousinweb}
\textbf{psallō} was borrowed from Greek \emph{psallo} \textquotedblleft{}to pluck the strings\textquotedblright{} (Lewis \& Short; EtymDB). English \emph{psalm} comes from the same Greek verb.
\end{cousinweb}

\begin{grammarlens}[title={What you've built}]
One more word for telling, explaining and describing at B1.
\end{grammarlens}

\subsection*{Your turn}
\begin{itemize}
\raggedright
  \item \emph{Say it:} \emph{psallō, psallere}
  \item \emph{Say it:} \emph{psallō, psallere}, once more
  \item \emph{Say it:} a sentence of your own with \emph{psallō, psallere}
  \item \emph{Recall:} say the Latin for a string, then the Latin for a lyre-player, then say \emph{psallō, psallere} again
\end{itemize}

\begin{tcolorbox}[breakable,colback=teal!4,colframe=teal!35!black,title={Before you move on}]
What does psallō, psallere mean? (\textquotedblleft{}To play the cithara, sing to strings\textquotedblright{}.) Say it once more.
\end{tcolorbox}
\section[(dialogue)]{Words from chapters 537–544 — a review}

\label{lesson:LA-R544-review-537-544}



\begin{quote}
Say the Latin for a quarrel and to reconcile.
\end{quote}

\subsection*{Your turn}
Say these aloud:

\begin{itemize}
\raggedright
  \item a quarrel, to reconcile, to give as a present, faithless, to abandon
  \item at leisure, to be free, on holiday, to relax, relaxation
  \item dice-playing, a dice-box, a piece in a board game, play, to wrestle
  \item a doll, a hoop, a spinning top, a hand-ball game, swimming
  \item a bath-keeper, the hot room, the warm room, the cold room, a pool
  \item a strigil, to anoint, a walk, to stay in the country, lovely
  \item singing, a song, an old song, to sing together, melodious
  \item a cithara, a singer to the lyre, a string, a lyre-player, to play the cithara
\end{itemize}

\begin{tcolorbox}[breakable,colback=teal!4,colframe=teal!35!black,title={Before you move on}]
A quarrel? (\textbf{iūrgium, iūrgiī}.) To reconcile? (\textbf{reconciliō, reconciliāre}.) To give as a present? (\textbf{dōnō, dōnāre}.) Faithless? (\textbf{īnfīdus, īnfīda, īnfīdum}.) To abandon? (\textbf{dēserō, dēserere}.) At leisure? (\textbf{ōtiōsus, ōtiōsa, ōtiōsum}.) To be free? (\textbf{vacō, vacāre}.) On holiday? (\textbf{fēriātus, fēriāta, fēriātum}.) To relax? (\textbf{relaxō, relaxāre}.) Relaxation? (\textbf{remissiō, remissiōnis}.) Dice-playing? (\textbf{ālea, āleae}.) A dice-box? (\textbf{fritillus, fritillī}.) A piece in a board game? (\textbf{latrunculus, latrunculī}.) Play? (\textbf{lūsus, lūsūs}.) To wrestle? (\textbf{luctor, luctārī}.) A doll? (\textbf{pūpa, pūpae}.) A hoop? (\textbf{trochus, trochī}.) A spinning top? (\textbf{turbō, turbinis}.) A hand-ball game? (\textbf{harpastum, harpastī}.) Swimming? (\textbf{natātiō, natātiōnis}.) A bath-keeper? (\textbf{balneātor, balneātōris}.) The hot room? (\textbf{caldārium, caldāriī}.) The warm room? (\textbf{tepidārium, tepidāriī}.) The cold room? (\textbf{frīgidārium, frīgidāriī}.) A pool? (\textbf{piscīna, piscīnae}.) A strigil? (\textbf{strigilis, strigilis}.) To anoint? (\textbf{unguō, unguere}.) A walk? (\textbf{ambulātiō, ambulātiōnis}.) To stay in the country? (\textbf{rūsticor, rūsticārī}.) Lovely? (\textbf{amoenus, amoena, amoenum}.) Singing? (\textbf{cantus, cantūs}.) A song? (\textbf{cantiō, cantiōnis}.) An old song? (\textbf{cantilēna, cantilēnae}.) To sing together? (\textbf{concinō, concinere}.) Melodious? (\textbf{canōrus, canōra, canōrum}.) A cithara? (\textbf{cithara, citharae}.) A singer to the lyre? (\textbf{citharoedus, citharoedī}.) A string? (\textbf{chorda, chordae}.) A lyre-player? (\textbf{fidicen, fidicinis}.) To play the cithara? (\textbf{psallō, psallere}.)
\end{tcolorbox}
